\documentclass[11pt,a4paper]{article}

\usepackage[utf8]{inputenc}
\usepackage[T1]{fontenc}
\usepackage[slovak]{babel}
\usepackage{lmodern}
\usepackage{geometry}
\usepackage{booktabs}
\usepackage{amsmath}
\usepackage{hyperref}
\usepackage{graphicx}
\usepackage{array}
\usepackage{listings}
\usepackage{xcolor}

\geometry{
left=2.5cm,
right=2.5cm,
top=2.5cm,
bottom=2.5cm
}

\lstset{
basicstyle=\ttfamily\small,
breaklines=true,
columns=fullflexible,
frame=single,
backgroundcolor=\color{gray!6}
}

\newcommand{\modelbase}{\texttt{baseline\_qwen25\_coder\_7b}}
\newcommand{\modeltest}{\texttt{b2\_testing\_v2}}
\newcommand{\modelref}{\texttt{b2\_refactoring\_v2}}
\newcommand{\modelshared}{\texttt{b1\_shared\_v2}}
\newcommand{\projectname}{\texttt{llm-ontology-experiment}}

\title{
\textbf{Zdieľané a špecializované LoRA adaptéry pre úlohy softvérového inžinierstva}\\
\large Report k predmetu Current Approaches in Machine Learning
}

\author{Patrik}
\date{\today}

\begin{document}

\maketitle

\begin{abstract}
Táto práca skúma, do akej miery dokáže menší open-source jazykový model pre kód zdieľať znalosti medzi dvoma úlohami softvérového inžinierstva: generovaním jednotkových testov a refaktoringom Java kódu. Experiment porovnáva základný model bez fine-tuningu, dva špecializované LoRA adaptéry a jeden spoločný LoRA adaptér trénovaný na kombinovaných dátach z oboch úloh. Špecializované modely boli trénované samostatne pre generovanie testov a refaktoring, zatiaľ čo spoločný model bol trénovaný na oboch doménach naraz.

Modely boli vyhodnocované pomocou automatických proxy metrík pre kvalitu testov, coverage proxy, code health, cohesion, coupling a refactoring quality. Okrem výkonnostnej evaluácie bola vykonaná aj analýza samotných LoRA adaptérov, ktorej cieľom bolo zistiť, či sa naučené aktualizácie váh líšia medzi úlohami a či je spoločný adaptér podobnejší testovaciemu alebo refaktoringovému adaptéru. Výsledky ukazujú, že špecializované adaptéry dosahujú najlepšie výsledky na svojich vlastných úlohách, zatiaľ čo spoločný model zostáva konkurencieschopný na oboch úlohách. Pri generovaní testov bola negatívna interferencia prakticky nulová, zatiaľ čo pri refaktoringu bola výhoda špecializovaného modelu výraznejšia.
\end{abstract}

\section{Úvod}

Veľké jazykové modely sa v posledných rokoch čoraz častejšie využívajú v softvérovom inžinierstve. Medzi typické úlohy patria generovanie kódu, generovanie jednotkových testov, refaktoring, oprava chýb, dokumentácia alebo vysvetľovanie programov. Plný fine-tuning veľkých modelov je však výpočtovo náročný a často nepraktický, najmä pri práci na bežne dostupnom hardvéri. Preto sú zaujímavé metódy parameter-efficient fine-tuning, medzi ktoré patrí aj LoRA a QLoRA.

Projekt \projectname{} bol vytvorený ako experimentálna pipeline pre diplomovú prácu zameranú na použitie jazykových modelov pri úlohách softvérového inžinierstva. Projekt pokrýva celý experimentálny tok: prípravu datasetov, instruction-tuning formát, LoRA/QLoRA tréning, inferenciu, výpočet metrík, reportovanie a následnú analýzu LoRA adaptérov.

Cieľom bolo preskúmať, či je pri dvoch príbuzných, ale odlišných úlohách výhodnejšie trénovať jeden spoločný model alebo viacero špecializovaných modelov. Konkrétne boli skúmané dve úlohy:
\begin{itemize}
\item generovanie jednotkových testov pre Java metódy,
\item refaktoring Java kódu.
\end{itemize}

Tieto úlohy majú spoločné to, že vyžadujú porozumenie zdrojovému kódu, metódam, identifikátorom a štruktúre programu. Zároveň sa však líšia typom požadovaného výstupu. Pri generovaní testov model vytvára nový JUnit testovací kód, zatiaľ čo pri refaktoringu model transformuje existujúci kód tak, aby zlepšil jeho štruktúru alebo zodpovedal určitému typu refaktoringu.

Hlavná výskumná otázka bola:

\begin{quote}
Do akej miery dokáže menší jazykový model pre kód zdieľať reprezentácie medzi generovaním testov a refaktoringom, a kedy je výhodnejšie použiť špecializované adaptéry?
\end{quote}

Na zodpovedanie tejto otázky boli porovnané štyri varianty modelu:
\begin{itemize}
\item základný model bez fine-tuningu,
\item špecializovaný LoRA adaptér pre generovanie testov,
\item špecializovaný LoRA adaptér pre refaktoring,
\item spoločný LoRA adaptér trénovaný na kombinovanom datasete.
\end{itemize}

\section{Projekt a experimentálna pipeline}

Repozitár \projectname{} je organizovaný tak, aby oddelil knižničnú logiku od spustiteľných skriptov. Hlavný kód sa nachádza v priečinku \texttt{src/llm\_ontology/}, zatiaľ čo priečinok \texttt{scripts/} obsahuje CLI vstupy pre prípravu dát, tréning a evaluáciu.

Dôležité časti projektu sú:
\begin{itemize}
\item \texttt{src/llm\_ontology/data/}: príprava datasetov Methods2Test, MaRV a ML4Refactoring,
\item \texttt{src/llm\_ontology/finetuning/}: prompt formatter, dataset loader a model loader,
\item \texttt{src/llm\_ontology/training/}: QLoRA tréningová logika,
\item \texttt{src/llm\_ontology/evaluation/}: inference, proxy metriky, reportovanie a analýzy,
\item \texttt{configs/}: YAML konfigurácie pre modely, tréning a evaluáciu.
\end{itemize}

Výsledkom projektu nie je iba jeden model, ale reprodukovateľná experimentálna infraštruktúra. Tá umožňuje pripraviť datasety, natrénovať LoRA adaptéry, vyhodnotiť ich na viacerých úlohách a analyzovať rozdiely medzi špecializovanými a zdieľanými adaptérmi.

\section{Datasety a formát úloh}

Experiment používa tri zdroje dát:
\begin{itemize}
\item Methods2Test pre generovanie JUnit testov,
\item MaRV pre refaktoringové príklady,
\item ML4Refactoring pre väčší refaktoringový dataset.
\end{itemize}

Spracované dáta sú uložené v instruction-tuning formáte. Každý záznam obsahuje inštrukciu, vstup, očakávaný výstup, doménu a zdroj. Zjednodušený príklad záznamu vyzerá nasledovne:

\begin{lstlisting}
{
  "instruction": "Vygeneruj JUnit test pre nasledujucu Java metodu.",
  "input": "public int add(int a, int b) { return a + b; }",
  "output": "@Test public void testAdd() { assertEquals(3, add(1, 2)); }",
  "domain": "testing",
  "source": "methods2test"
}
\end{lstlisting}

Pre refactoring má záznam podobný tvar, ale výstupom je upravená verzia kódu:

\begin{lstlisting}
{
  "instruction": "Vygeneruj refaktorovanu verziu nasledujuceho Java kodu podla typu refaktoringu: Rename Method.",
  "input": "public void oldName() { System.out.println(\"hello\"); }",
  "output": "public void newName() { System.out.println(\"hello\"); }",
  "domain": "refactoring",
  "source": "marv"
}
\end{lstlisting}

Finálne datasety použité pre tréning boli:
\begin{itemize}
\item \texttt{testing/}: Methods2Test pre B2-T model,
\item \texttt{refactoring/}: ML4Refactoring + MaRV pre B2-R model,
\item \texttt{combined/}: Methods2Test + ML4Refactoring pre B1 shared model.
\end{itemize}

\section{Promptovací formát a tréningová pipeline}

V2 experimenty boli vykonané po oprave tréningovej pipeline. Táto oprava bola dôležitá, pretože pilotné verzie modelov nepracovali s promptom a label maskovaním úplne konzistentne.

Použitý prompt formát bol jednotný pre tréning aj evaluáciu:

\begin{lstlisting}
### Instruction:
<instruction>

### Input:
<input code>

### Response:
<expected output>
\end{lstlisting}

Pri tréningu sa model nemá učiť reprodukovať zadanie. Má sa učiť generovať iba odpoveď. Preto sa tokeny patriace časti \texttt{Instruction} a \texttt{Input} v labeloch maskovali hodnotou $-100$. Strata sa tak počítala iba nad výstupnou časťou za \texttt{Response}.

Najdôležitejšie zmeny v2 pipeline boli:
\begin{itemize}
\item jednotný promptovací formát pre tréning aj evaluáciu,
\item použitie pôvodných inštrukcií z datasetu,
\item pridanie EOS tokenu na koniec cieľového výstupu,
\item výpočet straty iba nad odpoveďou modelu,
\item maskovanie prompt tokenov hodnotou $-100$,
\item bezpečné orezávanie dlhých príkladov tak, aby výstupné tokeny nezmizli.
\end{itemize}

Táto oprava bola overená pomocou debug skriptu, ktorý kontroloval počet prompt tokenov, počet full input tokenov, počet nemaskovaných label tokenov, EOS token a správanie \texttt{DataCollatorForSeq2Seq}. Pri najdlhšom refactoring príklade bolo overené, že output labels nezmizli celé.

\section{Základný model a LoRA fine-tuning}

Ako základný model bol použitý Qwen2.5-Coder-7B-Instruct dostupný lokálne ako:

\begin{lstlisting}
/mnt/c/models/huggingface/Qwen2.5-Coder-7B-Instruct
\end{lstlisting}

Namiesto plného fine-tuningu boli použité LoRA a QLoRA techniky, ktoré umožňujú adaptovať model pomocou malého počtu dodatočných parametrov.

Pri LoRA sa pôvodná váhová matica nemení priamo. Namiesto toho sa k nej pridáva nízkorozmerná aktualizácia:
\begin{equation}
W' = W + \Delta W,
\qquad
\Delta W = BA,
\end{equation}
kde $A$ a $B$ sú trénovateľné nízkorozmerné matice. Pôvodné váhy modelu zostávajú zmrazené a trénujú sa iba LoRA parametre.

Použité nastavenia LoRA/QLoRA boli:
\begin{itemize}
\item LoRA rank $r = 16$,
\item LoRA alpha $= 32$,
\item dropout $= 0.05$,
\item 4-bitová QLoRA kvantizácia typu NF4,
\item learning rate $2 \cdot 10^{-4}$,
\item early stopping podľa validačnej straty.
\end{itemize}

\section{Porovnávané modely}

Vo finálnej evaluácii boli porovnávané štyri modely uvedené v tabuľke~\ref{tab:models}.

\begin{table}[h]
\centering
\begin{tabular}{ll}
\toprule
Model & Popis \\
\midrule
\modelbase & Základný model bez LoRA adaptácie \\
\modeltest & LoRA adaptér trénovaný iba na generovaní testov \\
\modelref & LoRA adaptér trénovaný iba na refaktoringu \\
\modelshared & LoRA adaptér trénovaný na kombinovaných dátach \\
\bottomrule
\end{tabular}
\caption{Porovnávané modelové varianty.}
\label{tab:models}
\end{table}

Špecializované modely umožňujú zistiť, aký výkon možno dosiahnuť pri tréningu na jednej doméne. Spoločný model slúži na overenie, či jeden adaptér dokáže efektívne pokryť obe úlohy naraz.

\section{Evaluačné metriky}

\subsection{Metriky pre generovanie testov}

Pri generovaní testov boli použité proxy metriky, ktoré hodnotia štrukturálnu kvalitu vygenerovaných JUnit testov. Hodnotené boli najmä:
\begin{itemize}
\item prítomnosť anotácie \texttt{@Test},
\item prítomnosť assertionov,
\item volanie cieľovej metódy,
\item prítomnosť triviálnych testovacích vzorov,
\item priemerné test quality skóre,
\item priemerné coverage proxy skóre.
\end{itemize}

Coverage proxy nie je reálna line coverage ani branch coverage. Ide o automatický odhad, či vygenerovaný test svojou štruktúrou vyzerá ako test, ktorý by mohol overovať relevantné správanie.

\subsection{Metriky pre refaktoring}

Pri refaktoringu boli použité proxy metriky zamerané na textové a štrukturálne vlastnosti kódu. Hodnotené boli najmä:
\begin{itemize}
\item či sa výstup líši od vstupu,
\item podobnosť k referenčnému výstupu,
\item code health delta,
\item cohesion proxy,
\item coupling proxy,
\item priemerné refactoring quality skóre.
\end{itemize}

Tieto metriky umožňujú škálovateľné porovnanie modelov, ale nenahrádzajú plnohodnotnú statickú analýzu, kompiláciu ani overenie zachovania správania programu.

\section{Kvalitatívne príklady}

Nasledujúce príklady ilustrujú typ úloh, ktoré boli modelom predkladané. Prvý príklad reprezentuje generovanie testu:

\begin{lstlisting}
Instruction:
Vygeneruj JUnit test pre nasledujucu Java metodu.

Input:
public boolean isEven(int value) {
    return value % 2 == 0;
}

Expected output:
@Test
public void testIsEven() {
    assertTrue(isEven(2));
    assertFalse(isEven(3));
}
\end{lstlisting}

Druhý príklad reprezentuje refactoring:

\begin{lstlisting}
Instruction:
Vygeneruj refaktorovanu verziu nasledujuceho Java kodu podla typu refaktoringu: Extract Method.

Input:
public void printUser(User user) {
    System.out.println(user.getName());
    System.out.println(user.getEmail());
}

Expected output:
public void printUser(User user) {
    printUserDetails(user);
}

private void printUserDetails(User user) {
    System.out.println(user.getName());
    System.out.println(user.getEmail());
}
\end{lstlisting}

Tieto príklady sú zjednodušené, ale ukazujú rozdiel medzi úlohami. Pri testovaní model dopĺňa nový overovací kód. Pri refactoringu model mení štruktúru existujúceho kódu.

\section{Výsledky}

\subsection{Výsledky generovania testov}

Tabuľka~\ref{tab:testing-results} zobrazuje výsledky na 50 príkladoch z testovacej úlohy.

\begin{table}[h]
\centering
\small
\begin{tabular}{lrrrrrr}
\toprule
Model & @Test & Assert. & Target & Smell & Quality & Coverage \\
\midrule
\texttt{baseline} & 0.04 & 0.04 & 0.04 & 0.00 & 1.32 & 1.32 \\
\modeltest & 1.00 & 0.90 & 0.88 & 0.00 & 8.08 & 8.08 \\
\modelref & 0.78 & 0.72 & 0.70 & 0.00 & 6.90 & 6.90 \\
\modelshared & 1.00 & 0.80 & 0.86 & 0.00 & 8.06 & 8.06 \\
\bottomrule
\end{tabular}
\caption{Výsledky modelov na úlohe generovania testov. Hodnoty @Test, Assert., Target a Smell sú uvádzané ako podiely.}
\label{tab:testing-results}
\end{table}

Najlepšie výsledky dosiahol špecializovaný model \modeltest, ktorý získal priemerné test quality skóre 8.08. Spoločný model \modelshared dosiahol veľmi podobné skóre 8.06. Rozdiel medzi nimi je iba 0.02 bodu, čo naznačuje, že spoločný tréning na oboch doménach nespôsobil výrazné zhoršenie schopnosti generovať testy.

Refaktoringový model \modelref dosiahol na testovacej úlohe nižšie skóre 6.90. Napriek tomu stále produkoval čiastočne použiteľné testovacie výstupy. Základný model bez fine-tuningu dosiahol skóre iba 1.32, čo znamená, že v použitom promptovacom formáte spoľahlivo negeneroval očakávanú JUnit štruktúru.

\subsection{Výsledky refaktoringu}

Tabuľka~\ref{tab:refactoring-results} zobrazuje výsledky na 50 refaktoringových príkladoch.

\begin{table}[h]
\centering
\small
\begin{tabular}{lrrrrrr}
\toprule
Model & Differs & Edit sim. & Health & Cohesion & Coupling & Quality \\
\midrule
\texttt{baseline} & 1.00 & 0.01 & 6.06 & 6.10 & 0.20 & 4.31 \\
\modeltest & 0.48 & 0.68 & 5.54 & 9.15 & 7.78 & 5.96 \\
\modelref & 0.96 & 0.80 & 6.74 & 9.38 & 8.68 & 7.20 \\
\modelshared & 0.84 & 0.76 & 6.39 & 9.33 & 8.75 & 6.82 \\
\bottomrule
\end{tabular}
\caption{Výsledky modelov na úlohe refaktoringu. Differs označuje podiel výstupov, ktoré sa líšili od vstupného kódu.}
\label{tab:refactoring-results}
\end{table}

Na refaktoringovej úlohe dosiahol najlepší výsledok špecializovaný model \modelref s refactoring quality skóre 7.20. Výstup sa líšil od vstupu v 96~\% prípadov a model dosiahol najvyššiu podobnosť k referenčnému výstupu.

Spoločný model \modelshared dosiahol skóre 6.82, čo je stále konkurencieschopný výsledok, ale rozdiel oproti refaktoringovému špecialistovi je väčší než pri testovacej úlohe. Testing model \modeltest dosiahol na refactoringu skóre 5.96 a menil vstup iba v 48~\% prípadov. To naznačuje, že testovací adaptér je pri refactoringu konzervatívnejší a menej vhodný.

\section{Interferencia medzi úlohami}

Interferencia bola definovaná ako rozdiel medzi výkonom spoločného modelu a príslušného špecializovaného modelu:
\begin{equation}
I_{\text{task}} = S_{\text{shared}} - S_{\text{specialized}}.
\end{equation}
Negatívna hodnota znamená, že spoločný model zaostáva za špecializovaným modelom.

\begin{table}[h]
\centering
\begin{tabular}{lrrrr}
\toprule
Úloha & Shared & Specialized & Rozdiel & Relatívny rozdiel \\
\midrule
Testing & 8.06 & 8.08 & -0.02 & -0.25~\% \\
Refactoring & 6.82 & 7.20 & -0.38 & -5.23~\% \\
\bottomrule
\end{tabular}
\caption{Odhad interferencie spoločného modelu oproti špecializovaným modelom.}
\label{tab:interference}
\end{table}

Pri testovacej úlohe bola interferencia prakticky nulová. Spoločný model zaostal za špecializovaným testing modelom iba o 0.02 bodu, čo predstavuje relatívny pokles 0.25~\%. Pri refactoringu bol rozdiel väčší. Spoločný model zaostal za refaktoringovým špecialistom o 0.38 bodu, teda približne o 5.23~\%.

Z toho vyplýva, že shared fine-tuning je veľmi efektívny pri generovaní testov, no refactoring viac profituje zo špecializácie.

\begin{table}[h]
\centering
\begin{tabular}{lrr}
\toprule
Model & Testing score & Refactoring score \\
\midrule
\texttt{baseline} & 1.32 & 4.31 \\
\modeltest & 8.08 & 5.96 \\
\modelref & 6.90 & 7.20 \\
\modelshared & 8.06 & 6.82 \\
\bottomrule
\end{tabular}
\caption{Cross-task porovnanie modelov na oboch úlohách.}
\label{tab:cross-task}
\end{table}

Cross-task výsledky v tabuľke~\ref{tab:cross-task} ukazujú jasnú úlohovú špecializáciu. Testing model je najlepší na generovaní testov, ale slabší na refactoringu. Refactoring model je najlepší na refactoringu, ale nedosahuje výkon testing ani shared modelu na testovacej úlohe. Spoločný model je najvyrovnanejší.

\section{Analýza LoRA adaptérov}

Okrem výstupnej kvality bola vykonaná aj analýza samotných LoRA adaptérov. Cieľom bolo zistiť, ktoré vrstvy modelu boli jednotlivými úlohami najviac adaptované a či je spoločný adaptér podobnejší testing alebo refactoring adaptéru.

Analýza bola vykonaná priamo nad súbormi \texttt{adapter\_model.safetensors}. Pre každý LoRA modul boli extrahované matice \texttt{lora\_A} a \texttt{lora\_B}. Aktualizácia váh bola aproximovaná ako:
\begin{equation}
\Delta W = BA.
\end{equation}
Veľkosť zmeny bola vyhodnocovaná pomocou normy aktualizácie, prípadne proxy hodnoty:
\begin{equation}
\lVert A \rVert_F \cdot \lVert B \rVert_F.
\end{equation}
Podobnosť medzi adaptérmi bola meraná pomocou kosínusovej podobnosti zodpovedajúcich LoRA aktualizácií.

\subsection{Vrstvy s najväčšími LoRA zmenami}

Tabuľka~\ref{tab:lora-top-layers} zobrazuje päť vrstiev s najväčšou agregovanou LoRA normou pre každý adaptér.

\begin{table}[h]
\centering
\small
\begin{tabular}{llr}
\toprule
Model & Vrstva & LoRA norma \\
\midrule
\modeltest & 24 & 4.6321 \\
\modeltest & 27 & 4.5190 \\
\modeltest & 25 & 4.4920 \\
\modeltest & 23 & 4.4020 \\
\modeltest & 26 & 4.2982 \\
\midrule
\modelref & 25 & 15.9787 \\
\modelref & 21 & 15.6578 \\
\modelref & 20 & 15.0451 \\
\modelref & 22 & 14.6503 \\
\modelref & 19 & 14.1548 \\
\midrule
\modelshared & 23 & 10.4000 \\
\modelshared & 25 & 10.0997 \\
\modelshared & 24 & 10.0352 \\
\modelshared & 21 & 10.0150 \\
\modelshared & 22 & 9.8465 \\
\bottomrule
\end{tabular}
\caption{Vrstvy s najväčšou agregovanou LoRA normou pre jednotlivé adaptéry.}
\label{tab:lora-top-layers}
\end{table}

Testovací adaptér má najväčšie zmeny najmä vo vyšších vrstvách modelu, konkrétne vo vrstvách 23 až 27. Refaktoringový adaptér má najsilnejšie zmeny mierne nižšie, najmä vo vrstvách 19 až 25. Spoločný adaptér kombinuje oba vzory: prekrýva sa s testovacím adaptérom vo vrstvách 23, 24 a 25 a s refaktoringovým adaptérom vo vrstvách 21, 22 a 25.

\subsection{Podobnosť LoRA adaptérov}

Tabuľka~\ref{tab:lora-similarity} uvádza priemernú kosínusovú podobnosť medzi adaptérmi.

\begin{table}[h]
\centering
\small
\begin{tabular}{lr}
\toprule
Porovnávané adaptéry & Priemerná kosínusová podobnosť \\
\midrule
\modeltest{} vs. \modelref & 0.0096 \\
\modelshared{} vs. \modeltest & 0.0903 \\
\modelshared{} vs. \modelref & 0.0651 \\
\bottomrule
\end{tabular}
\caption{Priemerná kosínusová podobnosť medzi LoRA adaptérmi.}
\label{tab:lora-similarity}
\end{table}

Podobnosť medzi testovacím a refaktoringovým adaptérom je veľmi nízka, iba 0.0096. To naznačuje, že tieto dva špecializované adaptéry sa naučili výrazne odlišné aktualizácie váh. Spoločný adaptér je podobnejší obom špecializovaným adaptérom než sú špecializované adaptéry navzájom. Zároveň je mierne bližší testing adaptéru ako refactoring adaptéru.

Tento výsledok je konzistentný s výkonnostnou evaluáciou. Spoločný model takmer dorovnal testing špecialistu, ale pri refactoringu zaostal za refaktoringovým špecialistom výraznejšie.

\section{Diskusia}

Výsledky podporujú hypotézu, že generovanie testov a refaktoring zdieľajú určitú časť znalostí, no zároveň vyžadujú aj úlohovo špecifickú adaptáciu. Špecializované modely dosiahli najlepšie výsledky na svojich doménach. Spoločný model však dosiahol veľmi dobrý kompromis a najmä pri generovaní testov prakticky nezaostal za špecializovaným modelom.

Najdôležitejšie zistenia sú:
\begin{itemize}
\item špecializovaný testing adaptér dosiahol najlepšie výsledky na generovaní testov,
\item špecializovaný refactoring adaptér dosiahol najlepšie výsledky na refactoringu,
\item spoločný adaptér bol veľmi konkurencieschopný na oboch úlohách,
\item negatívna interferencia bola takmer nulová pri testingu a mierna pri refactoringu,
\item LoRA analýza ukázala veľmi nízku podobnosť medzi testing a refactoring adaptérom,
\item spoločný adaptér kombinuje adaptačné vzory oboch špecializovaných adaptérov.
\end{itemize}

Z praktického hľadiska to znamená, že spoločný adaptér môže byť efektívnym riešením v situáciách, kde je cieľom podporovať viacero príbuzných úloh naraz. Ak je však cieľom maximalizovať výkon na jednej konkrétnej úlohe, špecializovaný adaptér môže byť vhodnejší, najmä pri refactoringu.

\section{Limity práce}

Táto práca má niekoľko obmedzení. Po prvé, väčšina použitých metrík je proxy charakteru. Coverage proxy nenahrádza reálne meranie pokrytia pomocou JaCoCo. Podobne refactoringové metriky neoverujú sémantickú ekvivalenciu ani zachovanie správania programu.

Po druhé, finálne vyhodnotenie bolo vykonané na vzorke 50 príkladov pre každú úlohu. Táto veľkosť je vhodná pre projektové porovnanie a identifikáciu trendov, ale pre robustnejšie štatistické závery by bolo vhodné vyhodnotenie zopakovať na väčšej vzorke.

Po tretie, LoRA analýza bola vykonaná na úrovni váh adaptérov, nie na úrovni aktivácií modelu. Kosínusová podobnosť LoRA aktualizácií preto predstavuje iba aproximačný signál o podobnosti adaptérov a nie úplné mechanistické vysvetlenie správania modelu.

V ďalšej práci by bolo vhodné pripraviť menší executable subset, na ktorom by bolo možné vygenerované testy reálne kompilovať, spúšťať a merať pomocou JaCoCo. Pri refactoringu by bolo vhodné vytvoriť subset s existujúcimi testami, na ktorom by sa dalo overovať zachovanie správania pred a po refaktoringu.

\section{Záver}

V tejto práci boli porovnané zdieľané a špecializované LoRA adaptéry pre dve úlohy softvérového inžinierstva: generovanie jednotkových testov a refaktoring kódu. Výsledky ukázali, že špecializované adaptéry dosahujú najlepší výkon na svojich cieľových úlohách. Model \modeltest dosiahol najlepšie výsledky pri generovaní testov a model \modelref dosiahol najlepšie výsledky pri refactoringu.

Spoločný model \modelshared však dosiahol veľmi dobré výsledky na oboch úlohách. Pri generovaní testov bol jeho výkon prakticky rovnaký ako výkon špecializovaného testing modelu. Pri refactoringu síce zaostal za refaktoringovým špecialistom, ale stále bol druhým najlepším modelom.

Analýza LoRA adaptérov ukázala, že testovací a refaktoringový adaptér sa naučili veľmi odlišné aktualizácie váh. Spoločný adaptér sa nachádzal medzi nimi a kombinoval adaptačné vzory oboch úloh. Tieto výsledky naznačujú, že medzi test generation a refactoring existuje určitá miera zdieľateľných reprezentácií, ale maximálny výkon na konkrétnej úlohe stále profituje zo špecializácie.

\end{document}
